\begin{definition}
Let $n\ge0$ be an integer, let $H\colon\mathbb{C}\to\mathbb{C}$, let
$k\in\mathbb{Z}$ and let $r\in\mathbb{R}$.  Put
\[
  a_k(r)=\frac1{2\pi}\int_0^{2\pi}H\bigl(re^{i\theta}\bigr)\,e^{i(n-k)\theta}\,d\theta .
\]

This is the $k$-th Fourier coefficient (30) of the paper, of the function
$\theta\mapsto H(r,\theta)=H(re^{i\theta})$ relative to the system
$\bigl(e^{i(k-n)\theta}\bigr)_{k\in\mathbb{Z}}$ used in the expansion (29).
The scan of the article prints the upper limit of integration in (30) as
``$2k$''; the transcription records that it is $2\pi$, and that is the reading
used here.

The integral is the Lebesgue integral over the interval $[0,2\pi]$ and is read
as $0$ when the integrand is not integrable there; the paper uses $a_k(r)$ only
for $0\le r<1$ and for $H$ continuous on the unit disk, where the integrand is
continuous in $\theta$ and the integral is an ordinary one.  The shift by $n$ in
the exponent is the paper's: it is what makes $a_n$ the coefficient of the
constant mode, so that (23) $\lambda_n=\epsilon_n=0$ corresponds to the $k=n$
term of (19) being constant.
\end{definition}
